\section{Ristrutturazione dello Schema}

\subsection{Tabella dei volumi}

La tabella dei volumi costituisce lo strumento fondamentale 
per valutare i costi delle operazioni e analizzare le ridondanze, 
ponendo le basi per la corretta ristrutturazione dello schema E-R. 
Al suo interno viene stimata la quantità media di ciascun oggetto 
presente nella base di dati a metà del suo ciclo di vita operativo.

\begin{table}[h]
    \centering
    \renewcommand{\arraystretch}{1.2}
    \begin{tabular}{|l|c|r|}
        \hline
        \textbf{Concetto} & \textbf{Tipo} & \textbf{Volume} \\
        \hline
        Dipartimento & E & 8 \\
        \hline
        Corso di laurea & E & 80 \\
        \hline
        Docente & E & 700 \\
        \hline
        Insegnamento & E & 1.200 \\
        \hline
        Edizione & E & 2.400 \\
        \hline
        Lezione & E & 48.000 \\
        \hline
        Studente & E & 16.000 \\
        \hline
        Esame & E & 192.000 \\
        \hline
        è affiliato ad & R & 700 \\
        \hline
        è iscritto & R & 16.000 \\
        \hline
        nel piano di studio & R & 320.000 \\
        \hline
        sostiene & R & 192.000 \\
        \hline
        prevede & R & 192.000 \\
        \hline
        prerequisito & R & 360 \\
        \hline
        può insegnare & R & 3.500 \\
        \hline
        è erogato in & R & 2.400 \\
        \hline
        tiene & R & 2.800 \\
        \hline
        si compone di & R & 48.000 \\
        \hline
    \end{tabular}
    \caption{Tabella dei volumi}
    \label{tab:volumi}
\end{table}

La stima dei volumi (Tabella \ref{tab:volumi}) è stata ricavata da un'analisi quantitativa legata a una realtà accademica concreta, 
prendendo come modello di riferimento l'intero Ateneo dell'Università degli Studi di Udine.
I dati di riferimento sono stati estrapolati dagli Open Data ufficiali del Ministero dell'Università e della Ricerca (portale USTAT) 
relativi all'anno 2024 e dall'offerta formativa di ateneo per l'anno accademico 2023/2024. 

L'università ospita complessivamente circa 16.000 studenti, 8 dipartimenti distribuiscono l'offerta formativa su 80 corsi di laurea (tra triennali, magistrali e a ciclo unico) e 1.200 insegnamenti.
Il corpo docente è composto da circa 700 unità.

Partendo da queste basi fattuali, i volumi delle relazioni e delle entità intermedie sono stati ricavati matematicamente per un sistema a metà del suo ciclo di vita:
\begin{itemize}
    \item \textbf{Edizione e Lezione}: Mantenendo uno storico operativo di due anni accademici, i 1.200 insegnamenti generano 2.400 istanze per l'entità \textit{Edizione} e per la relazione \textit{è erogato in}. Assumendo una media di 20 appuntamenti per edizione, si ricavano $2400 \cdot 20 = 48000$ istanze per l'entità \textit{Lezione} e per la relazione \textit{si compone di}.
    \item \textbf{Piano di studio}: Si stima che uno studente inserisca in media 20 insegnamenti nel proprio piano di studi durante la sua carriera e considerando 16.000 studenti attivi, la relazione conta $320000$ tuple.
    \item \textbf{Esame, sostiene e prevede}: Considerando l'eterogeneità della popolazione studentesca, costituita da neo-immatricolati, studenti in pari e fuori corso, si stima una media ponderata di 12 esami sostenuti per studente. Di conseguenza, l'entità \textit{Esame} e le relazioni \textit{sostiene} e \textit{prevede} generano $192000$ istanze.
    \item \textbf{Può insegnare e tiene}: Ipotizzando che ogni docente sia abilitato a insegnare mediamente 5 insegnamenti afferenti al proprio settore disciplinare, la relazione \textit{può insegnare} genera $700 \cdot 5 = 3500$ tuple. Per la relazione \textit{tiene}, includendo le co-docenze, si stimano 2.800 assegnazioni sulle $2400$ edizioni attive.
    \item \textbf{Prerequisito}: Si ipotizza che all'interno di ciascuno degli 80 corsi di laurea vi siano in media 4 o 5 insegnamenti propedeutici, generando una stima di 360 istanze per la relazione ricorsiva.
\end{itemize}

\subsection{Analisi dei cicli}

Prima di procedere alla traduzione dello schema concettuale nel modello logico, 
viene effettuata una verifica strutturale del diagramma E-R per 
individuare eventuali percorsi ciclici. 
In particolare, sono stati individuati due cicli principali, 
successivamente analizzati per stabilire se costituissero ridondanze strutturali da eliminare 
oppure cicli giustificati da mantenere mediante opportuni vincoli di integrità nel modello logico.

\paragraph{Ciclo 1:}
Il primo ciclo si instaura tra le entità \textit{STUDENTE} - \textit{INSEGNAMENTO} - \textit{ESAME} - \textit{STUDENTE}. \\

\noindent
È possibile collegare uno studente a un insegnamento tramite due percorsi distinti:

\begin{enumerate}
    \item Il percorso diretto tramite la relazione \textit{nel piano di studio}.
    \item Il percorso indiretto tramite le relazioni \textit{sostiene} e \textit{prevede}.
\end{enumerate}

\noindent
Il ciclo non rappresenta una ridondanza, perché i due percorsi descrivono informazioni diverse.
Il primo percorso relativo ai corsi seguiti dallo studente ha uno scopo programmatorio, in quanto indica le attività che lo studente ha scelto nel proprio percorso di studi.
Il secondo percorso legato agli esami sostenuti, invece, ha uno scopo storico, poiché registra le prove realmente effettuate e concluse.
Eliminare uno dei due percorsi comporterebbe una perdita di informazioni, 
pertanto lo schema E-R viene mantenuto invariato. \\

\noindent
\textbf{Vincolo di integrità (Ciclo 1):} La correlazione logica tra i 
due percorsi impone che uno studente non possa sostenere l'esame di un 
corso non presente nel proprio piano di studi. Tale vincolo non è 
esprimibile tramite la struttura del modello E-R e viene pertanto 
demandato a un \textit{trigger}, che, prima di ogni inserimento nella 
relazione \textit{sostiene}, verifica che il corso associato 
all'esame risulti già presente nel piano di studi dello studente, 
in caso contrario l'operazione viene rifiutata.

\paragraph{Ciclo 2:}
Il secondo ciclo si instaura tra le entità \textit{DOCENTE} - \textit{INSEGNAMENTO} - \textit{EDIZIONE} - \mbox{\textit{DOCENTE}}. \\

\noindent
Un docente risulta collegato a un insegnamento tramite:

\begin{enumerate}
    \item Il percorso diretto tramite la relazione \textit{può insegnare}.
    \item Il percorso indiretto tramite le relazioni \textit{tiene} ed \textit{è erogato in}.
\end{enumerate}

\noindent
Anche in questo caso il ciclo non costituisce ridondanza strutturale. 
La relazione \textit{può insegnare} modella un'abilitazione statica, 
indicando le materie afferenti al settore scientifico-disciplinare del docente, 
indipendentemente dall'effettiva erogazione. 
Il percorso indiretto modella invece l'incarico reale assegnato in una precisa edizione. 
Essendo abilitazione e assegnazione due concetti temporalmente e funzionalmente separati, 
lo schema E-R viene mantenuto invariato. \\

\noindent
\textbf{Vincolo di integrità (Ciclo 2):} La correlazione logica tra i 
due percorsi impone che un docente non possa essere assegnato 
all'edizione di un corso per il quale non risulti abilitato. 
Anche questo vincolo viene demandato a un \textit{trigger}, che, prima di 
ogni inserimento nella relazione \textit{tiene},
verifica che il docente compaia nella relazione \textit{può insegnare} 
per il corso corrispondente all'edizione, altrimenti l'operazione viene rifiutata.

\subsection{Analisi delle ridondanze}

In questa fase della progettazione logica, si analizza l'impatto prestazionale 
derivante dal mantenimento degli attributi ridondanti 
del modello concettuale all'interno del modello logico. 
L'unico caso riguarda l'attributo derivato \textit{Crediti Ottenuti} in \textit{STUDENTE},  
il cui valore è calcolabile sommando i crediti degli insegnamenti 
per i quali lo studente ha registrato un esame superato (punteggio $\geq 18$). \\

\noindent
Per valutare se il risparmio computazionale in lettura 
giustifichi il sovraccarico in scrittura, 
l'analisi si basa su una stima delle frequenze di accesso ai dati.
Nello specifico, le operazioni coinvolte nell'analisi sono:
\begin{itemize}
    \item \textbf{Lettura:} verifica del totale dei crediti conseguiti da uno studente. 
    Stimando che ogni studente consulti il proprio stato 
    carriera mediamente 3 volte l'anno (inizio semestre, pre-sessione, verifica laurea), 
    si ottiene $16.000 \text{ studenti} \times 3 = 48.000 \approx 50.000$ esecuzioni/anno.
    \item \textbf{Scrittura:} aggiornamento dei crediti totali conseguiti da uno studente al superamento di un esame.  
    Stimando 5 esami superati all'anno per studente, si ottiene 
    $16.000 \text{ studenti} \times 5 = 80.000$ esecuzioni/anno.
\end{itemize}

\noindent
Per uniformare il confronto tra operazioni di natura diversa si adotta la convenzione $1w = 2r$ , 
ossia ogni accesso in scrittura equivale a due accessi in lettura, 
esprimendo così il costo in letture equivalenti.

\subsubsection{Accessi in assenza della ridondanza}

In assenza dell'attributo ridondante, per calcolare il totale dei 
crediti conseguiti da uno studente, il DBMS deve accedere a tutti 
gli esami superati in \textit{ESAME} (in media 12 record per studente) 
e, per ciascuno, recuperare il numero di crediti 
dall'entità \textit{INSEGNAMENTO}, complessivamente 
$12 \times 2 = 24$ accessi per esecuzione,
poiché ogni esame richiede una lettura in \textit{ESAME} 
e una in \textit{INSEGNAMENTO}.
$$24 \text{ letture} \times 50.000 \text{ volte/anno} = 1.200.000 \text{ accessi in lettura/anno}$$
L'inserimento di un nuovo esame richiede la sola scrittura 
di una tupla in \textit{ESAME}, ossia 2 accessi in lettura.
$$2 \text{ accessi} \times 80.000 \text{ volte/anno} = 160.000 \text{ accessi in lettura/anno}$$

\subsubsection{Accessi in presenza della ridondanza}

Mantenendo l'attributo ridondante, 
la verifica del totale crediti si riduce a 
un singolo accesso alla tupla dello studente, eliminando 
completamente la navigazione tra \textit{ESAME} e 
\textit{INSEGNAMENTO}.
$$1 \text{ lettura} \times 50.000 \text{ volte/anno} = 50.000 \text{ accessi in lettura/anno}$$
Al contrario, la scrittura si complica dato che oltre all'inserimento in \textit{ESAME} 
(1 scrittura), se il punteggio è $\geq 18$, un trigger deve leggere 
i crediti dell'insegnamento da \textit{INSEGNAMENTO} (1 lettura) e 
aggiornare i crediti dello studente (1 lettura e 1 scrittura), 
per un totale di 6 accessi considerando la conversione.
$$6 \text{ accessi} \times 80.000 = 480.000 \text{ accessi in lettura/anno}$$

\subsubsection{Decisione sulla ridondanza}

I costi complessivi annui espressi in accessi in lettura sono riassunti nella seguente tabella:

\begin{table}[h]
    \centering
    \renewcommand{\arraystretch}{1.2}
    \begin{tabular}{|l|r|r|}
        \hline
        \textbf{Operazione} & \textbf{Con ridondanza} 
        & \textbf{Senza ridondanza} \\
        \hline
        Lettura & 50.000 & 1.200.000 \\
        \hline
        Scrittura & 480.000 & 160.000 \\
        \hline
        \textbf{Totale} & \textbf{530.000} & \textbf{1.360.000} \\
        \hline
    \end{tabular}
    \caption{Bilancio accessi equivalenti per \textit{Crediti Ottenuti}}
\end{table}

\noindent
L'analisi evidenzia un vantaggio netto nel mantenere la ridondanza. 
L'abbattimento del 96\% dei costi di lettura (da 1.200.000 a 50.000 accessi/anno) 
compensa ampiamente il sovraccarico in scrittura (+320.000 accessi). 
Con un bilancio complessivo ridotto da 1.360.000 a 530.000 accessi annui,
viene scelto di mantenere \textit{Crediti Ottenuti} in \textit{STUDENTE}. 
La sua coerenza sarà garantita da un trigger 
attivato ad ogni inserimento in \textit{ESAME} con punteggio 
$\geq 18$, la cui implementazione è descritta nella sezione di 
progettazione fisica.
